% Responsable: @victoriaguilllen · Revisora: @gvidalquadras (issue #16)
\section{Introducción}\label{sec:intro}

Los transformers~\cite{vaswani2017attention} son la arquitectura dominante en
el aprendizaje profundo actual, y su pieza central es la capa de atención.
Dada una secuencia de $n$ tokens de dimensión $d$, apilados en una matriz
$X \in \R^{n\times d}$, la atención con softmax calcula
\begin{equation}\label{eq:atencion}
  A(X) = \softmax\!\left(\frac{X \Wq \Wk^{\top} X^{\top}}{\sqrt{d_k}}\right),
  \qquad X \mapsto A(X)\, X \Wv ,
\end{equation}
donde $\Wq, \Wk \in \R^{d\times d_k}$ y $\Wv \in \R^{d\times d}$ son matrices de
pesos y la softmax se aplica fila a fila. Por construcción, $A$ es
\emph{estocástica por filas}: sus entradas son positivas y $A\ones = \ones$,
de modo que cada token de salida es una media ponderada de los tokens de
entrada.

Que una red profunda se pueda entrenar depende en buena medida de cómo se
propagan la señal y los gradientes en la inicialización. En los transformers se
conocen dos patologías. La primera, común a muchas arquitecturas, es que los
gradientes se desvanezcan o exploten con la
profundidad~\cite{hanin2018which}. La segunda es el \emph{colapso de rango en
profundidad}: \textcite{dong2021attention} muestran que, en una red formada solo
por capas de atención, la salida converge de forma doblemente exponencial con
el número de capas a una matriz de rango~1, es decir, todos los tokens acaban
teniendo la misma representación. \textcite{noci2022signal} relacionan este
colapso con el desvanecimiento de los gradientes de las queries y las keys,
estudiando para ello el caso de atención uniforme.

\textcite{naitsaada2025mind} identifican un tercer fenómeno, propio de la atención
con softmax: el \emph{colapso de rango en anchura}. Incluso con una única capa,
las representaciones de los tokens se concentran en un único punto a medida que
crece la longitud de contexto~$n$. El origen es espectral. En la inicialización,
con pesos gaussianos y entrada isométrica ($XX^{\top} = I$), la matriz de
atención se comporta como una matriz de Markov aleatoria~\cite{bordenave2011circular}:
su mayor valor singular es~$1$, asociado a la dirección~$\ones$, mientras que
los demás son de orden $O(n^{-1/2})$. Este \emph{spectral gap} hace que
$A \approx \tfrac{1}{n}\ones\ones^{\top}$, una matriz que sustituye cada token
por la media de todos, y el rango estable~\cite{rudelson2007sampling} de la
salida tiende a~$1$. Visto así, suponer atención uniforme, como
en~\cite{noci2022signal}, equivale a partir de un colapso en anchura ya
consumado. Además, el gap agrava el colapso en profundidad y los problemas de
los gradientes. Los autores proponen un arreglo sencillo: eliminar el valor
atípico y usar $\Aperp = A - \tfrac{1}{n}\ones\ones^{\top}$, cuyo rango estable
crece linealmente con~$n$.

En este trabajo resumimos estos resultados y los reproducimos numéricamente. La
\cref{sec:gap} presenta las definiciones y el resultado principal sobre el gap
espectral y el colapso en anchura. La \cref{sec:profundidad} estudia sus
consecuencias en profundidad y en los gradientes, y la \cref{sec:arreglo}, el
arreglo propuesto. Cerramos con conclusiones y limitaciones en la
\cref{sec:conclusiones}. A lo largo del informe, $n$ denota la longitud de
contexto, que en~\cite{naitsaada2025mind} se llama~$T$. El código de los
experimentos está en el repositorio del proyecto.
